% HDBSCAN in five steps on 12 invented points: groups A and B (5 points each), two isolated
% points n1, n2. Adapted from earlier figures by the same author.
\begin{tikzpicture}[fig]
\def\W{2.5}\def\H{1.9}
\newcommand{\frm}{\draw[black!30] (0,0) rectangle (\W,\H);}
\newcommand{\pcap}[1]{\node[small, text width=2.75cm, anchor=north] at ({\W/2},-0.08) {#1};}
% the points: A, B, n1, n2
\newcommand{\Apts}{0.35/1.2, 0.55/1.42, 0.65/1.1, 0.45/0.95, 0.8/1.3}
\newcommand{\Bpts}{1.75/0.6, 1.95/0.8, 2.05/0.5, 1.8/0.38, 2.2/0.7}
\newcommand{\pts}[1]{% #1 = 1: colour the final groups
  \foreach \x/\y in \Apts {\ifnum#1=1 \fill[cK1] (\x,\y) circle (1.4pt); \else \fill[black!55] (\x,\y) circle (1.4pt); \fi}
  \foreach \x/\y in \Bpts {\ifnum#1=1 \fill[cK2] (\x,\y) circle (1.4pt); \else \fill[black!55] (\x,\y) circle (1.4pt); \fi}
  \ifnum#1=1
    \foreach \x/\y in {1.2/1.62, 2.25/1.55} \draw[black!60, thick] (\x,\y) ++(-0.06,-0.06) -- ++(0.12,0.12) ++(-0.12,0) -- ++(0.12,-0.12);
  \else
    \foreach \x/\y in {1.2/1.62, 2.25/1.55} \fill[black!55] (\x,\y) circle (1.4pt);
  \fi}
% ---------- 1. core distances ----------
\begin{scope}
  \frm \pts{0}
  \draw[cK1!80!black, thick] (0.65,1.1) circle (0.26);
  \draw[cAcc, thick, densely dashed] (1.2,1.62) circle (0.55);
  \node[small, cAcc, anchor=north] at (1.2,1.05) {large};
  \node[lab] at ({\W/2},2.2) {1. core distance};
  \pcap{distance to the $m$-th neighbour ($m = 3$): small in a crowd}
\end{scope}
% ---------- 2. minimum spanning tree ----------
\begin{scope}[xshift=3.05cm]
  \frm
  \draw[black!60] (0.35,1.2) -- (0.45,0.95) -- (0.65,1.1) -- (0.8,1.3) -- (0.55,1.42);
  \draw[black!60] (1.8,0.38) -- (1.75,0.6) -- (1.95,0.8) -- (2.2,0.7) -- (2.05,0.5);
  \draw[cAcc, thick] (0.8,1.3) -- (1.2,1.62);
  \draw[cAcc, thick] (1.95,0.8) -- (2.25,1.55);
  \draw[cAcc, very thick] (0.65,1.1) -- (1.75,0.6);
  \pts{0}
  \node[lab] at ({\W/2},2.2) {2. spanning tree};
  \pcap{tree of $d_{\rm mr}$; long edges (orange) join sparse parts}
\end{scope}
% ---------- 3. single-linkage hierarchy ----------
\begin{scope}[xshift=6.1cm]
  \draw[->] (0,0) -- (0,1.95) node[below left, small] {$d_{\rm mr}$};
  \draw[black!30] (0,0) -- (\W,0);
  \draw[cK1!80!black] (0.15,0) |- (0.35,0.33) -- (0.35,0) ;
  \draw[cK1!80!black] (0.55,0) |- (0.75,0.37) -- (0.75,0);
  \draw[cK1!80!black] (0.25,0.33) |- (0.65,0.45) -- (0.65,0.37);
  \draw[cK1!80!black] (0.45,0.45) |- (0.95,0.51) -- (0.95,0);
  \draw[black!60] (0.70,0.51) |- (1.15,0.99) -- (1.15,0);
  \draw[cK2!80!black] (1.35,0) |- (1.55,0.32) -- (1.55,0);
  \draw[cK2!80!black] (1.75,0) |- (1.95,0.35) -- (1.95,0);
  \draw[cK2!80!black] (1.45,0.32) |- (1.85,0.44) -- (1.85,0.35);
  \draw[cK2!80!black] (1.65,0.44) |- (2.15,0.53) -- (2.15,0);
  \draw[black!60] (1.90,0.53) |- (2.35,1.11) -- (2.35,0);
  \draw[cAcc, thick] (0.925,0.99) |- (2.125,1.71) -- (2.125,1.11);
  \node[small, cK1!80!black] at (0.55,-0.14) {A};
  \node[small] at (1.15,-0.14) {$n_1$};
  \node[small, cK2!80!black] at (1.75,-0.14) {B};
  \node[small] at (2.35,-0.14) {$n_2$};
  \node[lab] at ({\W/2},2.2) {3. hierarchy};
  \begin{scope}[yshift=-0.12cm]
  \pcap{cut edges from the longest: groups split as $d_{\rm mr}$ is lowered}
  \end{scope}
\end{scope}
% ---------- 4. condensed tree ----------
\begin{scope}[xshift=9.15cm]
  \draw[->] (0,1.9) -- (0,0) node[below left, small] {$\lambda$};
  \node[small, anchor=east] at (0,1.75) {0};
  \fill[black!20] (0.15,1.85) rectangle (2.35,1.45);
  \fill[cK1!40] (0.15,1.45) -- (1.15,1.45) -- (1.15,1.05) -- (0.97,1.05) -- (0.75,0.3)
    -- (0.55,0.3) -- (0.33,1.05) -- (0.15,1.05) -- cycle;
  \fill[cK2!40] (1.35,1.45) -- (2.35,1.45) -- (2.35,0.95) -- (2.17,0.95) -- (1.95,0.25)
    -- (1.75,0.25) -- (1.53,0.95) -- (1.35,0.95) -- cycle;
  \draw[black!60] (0.15,1.45) -- (2.35,1.45);
  \node[small] at (1.25,1.65) {all 12};
  \draw[->, black!60] (1.15,1.05) -- ++(0.25,-0.25) node[below, small] {$n_1$ out};
  \node[small, cK1!80!black] at (0.65,0.85) {A};
  \node[small, cK2!80!black] at (1.85,0.75) {B};
  \node[lab] at ({\W/2},2.2) {4. condensed tree};
  \pcap{$\lambda = 1/d_{\rm mr}$; width: points still in; tiny splits: points falling out}
\end{scope}
% ---------- 5. selected groups ----------
\begin{scope}[xshift=12.2cm]
  \frm \pts{1}
  \node[lab] at ({\W/2},2.2) {5. most stable};
  \pcap{branches of largest area (stability) kept; rest: noise ($\times$)}
\end{scope}
% ---------- formulas ----------
\node[note, text width=15cm, align=center, anchor=north] at (7.35,-1.55)
  {$\mathrm{core}_m(a)$ = distance to the $m$-th nearest neighbour, \qquad
   $d_{\rm mr}(a,b) = \max\big(\mathrm{core}_m(a), \mathrm{core}_m(b), d(a,b)\big)$,\\[2pt]
   stability$(C) = \sum_{p \in C} \big(\lambda_p - \lambda_{\rm birth}(C)\big)$, with $\lambda_p$
   the level at which point $p$ leaves the group};
\end{tikzpicture}
